\chapter{Conclusion and Future Work}
\label{ch:conclusion}

This thesis built, validated, and optimised a from-scratch CUDA/OptiX volumetric path tracer for
Gaussian kernel-mixture volumes---a high-performance re-implementation of \emph{Don't Splat Your
Gaussians}~\cite{DSYG}---one that, at equal quality, renders faster than the Mitsuba reference.

\section{Summary of contributions}
\label{sec:contributions}

The central contribution is a rendering architecture (\Cref{ch:architecture}) in two
parts. The first is \emph{single-trace any-hit collection}: it gathers all of a
ray's primitive entries in one acceleration-structure traversal, back-face-culled to entries, with
exits computed analytically. The second is \emph{analog-decomposition scatter sampling}, suggested by Condor, the lead author of the reference (personal communication, 2026). Each primitive draws a
closed-form free flight, and the nearest---the argmin---is the scattering event. Together they replace
the segment-marching and root-finding that the reference uses for the scatter
distance (validated end-to-end and by complexity argument rather than by internal ablation;
\Cref{ch:optimization} states this scope explicitly). Two further contributions support it. The
first is a \emph{differential-validation methodology}
(\Cref{ch:validation}) that establishes correctness against Mitsuba and, where a closed form exists,
against analytic ground truth. It doubles as the unbiasedness evidence for the new sampler. The second is the \emph{performance-engineering study} of \Cref{ch:optimization}. That study
contributed a second, scene-dependent win (volumetric product-RIS direct lighting,
$1.48\times$ at equal quality under environment lighting, default off) and a characterised
set of negative results.

At equal
quality the renderer is \num{2.7}$\times$ faster than the reference's corrected fast estimator on the
showcase, and it renders firefly-free at production sample budgets. That like-for-like matchup is
enabled by this work's reference corrections, submitted upstream. The advantage is not confined to
frame time. At equal conditions the shipped reference's faster direct-lighting estimator does not
converge to the correct image on dense media. This renderer's is both correct and low-variance
(\Cref{sec:results-firefly}). Per-asset cap calibration puts the renderer's peak
device memory below the reference's on the showcase (\Cref{sec:results-memory}). And its
ahead-of-time-compiled pipeline starts in roughly half the reference's warm start
(\Cref{sec:results-startup}). Finally, a cross-architecture probe (\Cref{sec:ser}) shows the
hand-written megakernel exploits Shader Execution Reordering on Ada hardware for a further
\num{1.12}--\num{1.68}$\times$ at equal quality, image-identical---a reordering the reference's compiled pipeline cannot request. Measured end-to-end on that hardware, the renderer's
configurations take the \num{64}-spp showcase
from \SI{3.29}{\second} (MIS) to \SI{2.31}{\second} with reordering and \SI{1.63}{\second} in the
maximal-throughput configuration (\Cref{sec:ser}).

The upstream patch set is itself part of the contribution: five changes, each backed by
deterministic evidence, submitted as a pull request~\cite{KramarzVolprimFixes2026}
(\Cref{tab:revision-timeline} gives the timeline). After the corrections, the two renderers agree
at the Monte-Carlo noise level (\Cref{fig:agreement-panels}). The headline like-for-like figure (\Cref{tab:likeforlike}) depends on that certification.
The renderer itself is also a contribution: an object-oriented codebase, RAII
throughout, whose experimental features sit behind guards that leave the production
binary byte-identical. Every experiment is reproducible from a documented invocation, and the
campaign scripts are kept alongside the sources (\Cref{sec:gpu-impl}). The codebase is released as
open source together with its test scenes.

\section{Limitations}
\label{sec:limitations}

Several limitations bound the renderer's current scope:

\begin{itemize}
  \item \textbf{No emission.} It models absorption and scattering but not emissive media, so combustion
  assets render as dark plumes; the assets carry no emission data either.
  \item \textbf{Gaussian kernel only.} The \emph{forward} optical depth is closed-form for other kernel
  families too (e.g.\ the compactly-supported Epanechnikov), but the argmin sampler needs a practical
  closed-form \emph{inverse} (\Cref{sec:analytic-od}), which is Gaussian-specific---for the Epanechnikov
  the inverse is impractical (the reference uses a numerical solver even for a single such kernel)---so
  other kernels are unsupported here; the architecture could accommodate one given its inverse, but does
  not implement it. Relatedly, assets fit with the Gabor-noise extension (Gabor kernels: Gaussians modulated by a cosine
carrier, adding high-frequency detail) render only their coarse
  Gaussian base; high-frequency Gabor detail is unsupported.
  \item \textbf{Static, single-frame.} There is no animation support.
  \item \textbf{Compile-time, asset-dependent caps.} The per-ray buffers (\Cref{sec:gpu-impl}) are sized
  at compile time; a substantially denser scene requires re-sizing and a rebuild. An offline estimator
  now predicts the required caps from an asset's geometry---the \num{25600}-Gaussian bunny, for one,
  needs markedly larger caps than the \num{652}-primitive cloud---so the sizing is determined by
  measurement rather than assumption, but the caps remain static, and the runtime overflow counter only
  \emph{detects} an undersized cap rather than fixing it.
  \item \textbf{Ampere hardware.} The one lever that directly targets the renderer's ray divergence,
  OptiX Shader Execution Reordering~\cite{NvidiaSER2022}, is Ada-only and unavailable on the RTX~3090
  used throughout---though its effect is measured directly on an Ada GPU in \Cref{sec:ser}
  (\num{1.12}--\num{1.68}$\times$, image-identical).
  \item \textbf{Small systematic residuals.} A handful of sub-$0.2\%$ systematic differences from the
  reference are documented rather than resolved, because each remains unattributed. The largest
  is a ${+}2\times10^{-4}$ difference at heavily-overlapped scatter vertices, measured
  \emph{not} to be a single-precision $\operatorname{erf}^{-1}$ artefact (a double-precision
  inverse moves it by ${\sim}2\times10^{-8}$). Its side is unknown: corrections were later made
  on both sides of the comparison in exactly this regime (\Cref{sec:results-firefly}), and
  \Cref{ch:validation} names the candidate causes. Also documented: a ${+}1.8\times10^{-4}$
  brightness systematic in low-density interiors, and a small channel-dependent residual in
  coloured-albedo rendering (cause not isolated). Each is bounded well under a part in
  $500$ and confined to a specific regime. Resolution has a known path---extending the exact
  float64 predictor of \Cref{app:corrections} from absorption to scattering on an overlap
  extract, and per-channel bisection of the coloured-albedo path---but neither was undertaken
  here.
\end{itemize}

\section{Future work}
\label{sec:future-work}

The most direct continuations are natural extensions of the kernel-mixture representation, several
suggested by Condor (personal communication): a \emph{unified emissive/absorptive kernel} model that exploits closed-form
integrals for accurate product sampling of light sources and BSDFs; \emph{SGGX and microflake} models
to represent surface-like materials within the same primitive framework; \emph{parametric kernels} for more compact and accurate radiance fields (a direction
the DSYG paper's own future work names~\cite{DSYG}); and a bridge from \emph{Lagrangian fluid simulation} directly
to kernel mixtures, bypassing voxelisation and kernel regression.

Three engineering directions follow from this thesis specifically. First, \emph{differentiable and
inverse rendering}: re-expressing the renderer with analytic backward passes or in
SLANG.D~\cite{SLANG.D}, to recover scene parameters from images---the inverse-rendering capability the
original proposal targeted. Second, \emph{productionising Shader Execution Reordering on Ada}:
\Cref{sec:ser} already measures a \num{1.12}--\num{1.68}$\times$ image-identical speedup from a minimal
per-bounce reorder on a rented Ada GPU; integrating it as a hardware-detected default path with a
per-scene-tuned coherence hint is the natural follow-through for Ada-class deployment. Third, \emph{recompile-free capacity}: the per-ray buffers are compile-time constants whose size the
offline estimator already predicts per asset (\Cref{tab:overlap}). Removing the fixed buffer
\emph{entirely} was tried and abandoned---the cap-free streaming variant (\Cref{sec:autopsies}) was bit-exact but slower. An untried middle path keeps a \emph{small, fixed} buffer but, when it fills,
re-issues the traversal from the last processed distance to gather the next chunk and repeats: this pays
an extra re-traversal only on the rare overflowing ray, keeps the single-trace fast path for the common
case, and removes the per-asset recompile without re-introducing the cap-free design's
whole-render cost. The bounce-\num{0} containment scan of \Cref{sec:results-scaling} has already yielded:
precomputing the camera's containing set once per render, implemented post-freeze on a
branch, is byte-identical in output and measured $1.1$--$3.0\times$ per asset at the
matched configuration ($2.95\times$ on the bunny). Merging it and re-baselining the
timing results, plus the general acceleration-structure point query for arbitrary-origin
cameras (the orthographic showcase still scans), is the natural first step. Path guiding, deferred here for want of a cheap
predictive test (\Cref{sec:autopsies}), would also merit a revisit with a full offline oracle.

\section{Closing}
\label{sec:closing}

The renderer demonstrates that ray-traced Gaussian kernel-mixture volumes can be rendered
both \emph{correctly} and \emph{efficiently} by exploiting the analytic, invertible
per-primitive optical depth: the resulting path tracer needs no sorting, no marching, and
no numerical root-finding. The negative results of \Cref{ch:optimization} are as useful in
practice: this workload is latency-bound, belongs in a single persistent kernel, and
resists the restructurings that help more conventional path tracers---which should spare
the next implementer the same experiments. The renderer, its test scenes, and the scripts behind
every figure in this document are available as open source (\Cref{sec:gpu-impl}), so the
results can be re-run, extended, or challenged.
